\documentclass[10pt,a4paper]{amsart}
\usepackage[margin=19mm]{geometry}
\usepackage{amsmath,amssymb,mathtools}
\usepackage[scr=rsfs,cal=euler]{mathalfa}
\usepackage{xfrac}
\usepackage[unicode,hidelinks,bookmarks=false]{hyperref}
\newcommand{\ie}{i.e.\ }
\newcommand{\cf}{cf.\ }
\newcommand{\resp}{respectively\ }
\newcommand{\Subsection}{\subsection}
\providecommand{\nobreakdash}{\nobreak\hbox{-}}
\setlength{\emergencystretch}{2em}
\multlinegap0pt
\usepackage{bm}
\usepackage{bbm}
\usepackage{dsfont}
\usepackage{xspace}
\usepackage{cancel}
\usepackage{mathtools}
\usepackage{amsthm}
\makeatletter
\@ifundefined{theorem}{\newtheorem{theorem}{Theorem}}{}
\@ifundefined{proposition}{\newtheorem{proposition}[theorem]{Proposition}}{}
\makeatother
\def\jp#1{{\left\langle{#1}\right\rangle}}
\renewcommand{\Im}{\textrm{Im}}
\title[Global hypoellipticity and global solvability of Vekua-type ]{Global hypoellipticity and global solvability of Vekua-type operators associated with diagonal operators on compact Lie groups}
\author{Ricardo Paleari da Silva}
\date{Source: arXiv:2604.06544v1 (CC BY 4.0)}
\begin{document}
\maketitle

\noindent\emph{Source and licence.} Global hypoellipticity and global solvability of Vekua-type operators associated with diagonal operators on compact Lie groups. Version arXiv:2604.06544v1 (CC BY 4.0), distributed under the Creative Commons Attribution 4.0 International licence, as recorded in the arXiv licence metadata for this exact version and shown on the abstract page https://arxiv.org/abs/2604.06544v1. Full rights notice: licenses/arxiv-2604.06544-CC-BY-4.0.txt.\par\smallskip

\noindent\textit{Editorial scope.} Counted unit: the Proposition whose source label is \texttt{None} in the article (source0.tex lines 327--328, source label \texttt{None}), delivered together with the article's own complete original proof of it (source0.tex lines 330--374, one \texttt{proof} environment of 45 lines). No mathematics is written, repaired or re-derived: statement and proof are the article's own text line for line. Required context retained uncounted: subsystem (lines 237--239); Zinfinite (lines 294--295). Every definition, display and label of those lines is retained unchanged. Omitted from this package and not redistributed: 1--236, 240--293, 296--326, 329--329, 375--900. Nothing inside a retained excerpt is omitted or altered: every retained body is delivered exactly as the article's lines are written, so no omission receipt of any kind accompanies this package. Citations: the retained excerpts carry no citation command at all, so no bibliography entry of the article is transcribed and no reference is redistributed. Reference adaptation: none was needed and none was made; every reference inside the delivered unit resolves inside it, so no unresolved reference or citation is produced. Printed slips kept literal: no printed word, symbol, hypothesis or number of the counted statement or of its proof was altered. Layout and class: the article's own class is \texttt{sn-jnl} with packages including indentfirst, enumerate, amssymb, amsfonts, amsmath, amsthm, mathrsfs, dsfont, color, graphicx, xcolor, comment, verbatim; the wrapper is the lane's pinned \texttt{amsart} editorial wrapper at 10pt on A4, which restates the theorem-like environments the retained text needs and adds the article's own macro definitions that the retained text needs (\texttt{\textbackslash Im}, \texttt{\textbackslash jp}), with extra wrapper packages (bm, bbm, dsfont, xspace, cancel, mathtools). The delivered numbering is the unit's own. Assets: no retained span contains a figure environment, a graphics-inclusion command, a PDF-page inclusion, a table or a TikZ picture; no image file is copied into this package, the package declares no asset dependency and no image file is redistributed with it. Licence: version arXiv:2604.06544v1 of this article is distributed under the Creative Commons Attribution 4.0 International licence, as recorded in the arXiv licence metadata for this exact version and shown on the abstract page https://arxiv.org/abs/2604.06544v1. A full rights notice is delivered at licenses/arxiv-2604.06544-CC-BY-4.0.txt. Characters: every retained excerpt is pure ASCII, so the delivered engine, XeLaTeX with its default Latin Modern text font, reports no missing character.
\begin{equation}\label{subsystem}
	\left \{ \begin{array}{c} (\sigma_k(\xi)-q)\cdot \widehat{u}(\xi)_{k\ell}-p \cdot \widehat{\overline{u}}(\xi)_{k\ell}=\widehat{f}(\xi)_{k\ell} \\ -\overline{p} \cdot \widehat{u}(\xi)_{k\ell} + (\sigma_k(\xi)-\overline{q} ) \cdot \widehat{ \overline{u}}(\xi)_{k\ell} = \widehat{\overline{f}}(\xi)_{k\ell}   \end{array}   \right.
\end{equation}

\begin{proposition}\label{Zinfinite} Suppose that $[\xi] \neq [\overline{\xi}]$ for all $\xi \in \textrm{Rep}(G)$. If the set $\mathcal{Z}$ is infinite, then $P$ is not (GH).	
\end{proposition}  

\begin{proposition} Suppose that $[\xi] \neq [\overline{\xi}]$ for all $\xi \in \textrm{Rep}(G)$. If the Vekua-type operator $P$ is (GH), then the condition (DC) holds.
\end{proposition} 

\begin{proof} Suppose that condition (DC) does not hold. Then, for each $n \in \mathbb{N}$ there exists $[\xi_n] \in \widehat{G}$ and $1 \leq k_n \leq d_{\xi_n}$ such that $\jp{\xi_n}>n$ and $|\Delta_{k_n}(\xi_n)|<\jp{\xi_n}^{-n}$ for all $n \in \mathbb{N}$. By Proposition \ref{Zinfinite} we may also assume that $|\Delta_{k_n}(\xi_n)|>0$ for all $n \in \mathbb{N}$. Now consider the sequence of matrices $(x(\xi))_{[\xi] \in \widehat{G}}$ given by
	$$x(\xi)_{k\ell}:= \left\{ \begin{array}{ccc} \Delta(\xi)_{k} & \textrm{if} & [\xi]=[\xi_n] \textrm{ or } [\xi]=[\overline{\xi}_n] \textrm{ and } k=\ell=k_n\\ 0 & & \textrm{otherwise}   \end{array}  \right.$$
By hypothesis it is clear that $(x(\xi))_{[\xi]}$ is the sequence of Fourier-coefficients of a smooth function $f \in C^\infty(G)$. Now, for each $c \in \mathbb{C}\setminus \{0\}$ consider the smooth function $f_c:=c \cdot f$ and the sequence of matrices $(y(\xi))_{[\xi]}$ given by
$$y(\xi)_{k\ell}:= \left\{ \begin{array}{ccc} \sigma_{k}(\xi)\cdot c-\overline{q}\cdot c+p\cdot \overline{c} & \textrm{if} & [\xi]=[\xi_n] \textrm{ or } [\xi]=[\overline{\xi}_n] \textrm{ and } k=\ell=k_n\\ 0 & & \textrm{otherwise}   \end{array}  \right.$$
Notice that for $[\xi]=[\xi_n]$ or $[\xi]=[\overline{\xi_n}]$ and $k=\ell=k_n$ we have
$$|y(\xi)_{k \ell}| \leq |c| \cdot(|\sigma_k(\xi)|+|q|+|p|) \leq |c| \cdot (C_\sigma +|p|+|q|)\cdot \jp{\xi}^m,$$
which implies that $(y(\xi))_{[\xi]}$ is the sequence of Fourier-coefficients of a distribution $u_c \in \mathcal{D}'(G)$. We claim that $Pu_c=f_c$. To see this, it is sufficient to verify that $\widehat{Pu_c}(\xi_n)_{k_n k_n}=\widehat{f_c}(\xi_n)_{k_n k_n}$ and $\widehat{Pu_c}(\overline{\xi_n})_{k_n k_n}=\widehat{f_c}(\overline{\xi_n})_{k_n k_n}$  for all $n$, $k,\ell$. But, again by equations (\ref{subsystem}) we have
\begin{eqnarray*}
	\widehat{Pu_c}(\xi_n)_{k_n k_n} & = & (\sigma_{k_n}(\xi_n)-q) \cdot \widehat{u_c}(\xi_n)_{k_n k_n}-p\cdot \widehat{\overline{u_c}}(\xi_n)_{k_n k_n} \\
	& = & (\sigma_{k_n}(\xi_n)-q) \cdot (\sigma_{k_n}(\xi_n)\cdot c-\overline{q}\cdot c +p\cdot \overline{c})- p \cdot \overline{(\sigma_{k_n}(\overline{\xi_n})\cdot c-\overline{q}\cdot c+p\cdot \overline{c})} \\
	& = & (\sigma_{k_n}(\xi_n)-q) \cdot (\sigma_{k_n}(\xi_n)\cdot c-\overline{q}\cdot c +p\cdot \overline{c})-p \cdot (\sigma_{k_n}(\xi_n)\cdot \overline{c}-q \cdot \overline{c}+\overline{p}\cdot c) \\
	& = & c\cdot (\sigma_{k_n}(\xi_n)-q)\cdot (\sigma_{k_n}(\xi_n)-\overline{q})+p \cdot \overline{c}\cdot (\sigma_{k_n}(\xi_n)-q)- p \cdot \overline{c} \cdot(\sigma_{k_n}(\xi_n)-q) \\ 
	&   & -c \cdot |p|^2 \\
	& = & c \cdot \Delta_{k_n}(\xi_n) \\
	& = & \widehat{f_c}(\xi_n)_{k_n k_n}
\end{eqnarray*} 
for all $n \in \mathbb{N}$. Similarly, it can be verified that $\widehat{Pu_c}(\overline{\xi_n})_{k_n k_n}=\widehat{f_c}(\overline{\xi_n})_{k_n k_n}$, concluding that $Pu_c=f_c$. 
Notice that if $u_c \in C^\infty(G)$, in particular we must have $\lim_n |\widehat{u_c}(\xi_n)_{k_n}| = \lim_n |\widehat{u_c}(\overline{\xi_n})_{k_n}|=0$. But since
$$|\widehat{u}(\xi_n)_{k_n}| = |c| \cdot \left|\sigma_{k_n}(\xi_n)-\left(\overline{q}-\frac{p \cdot \overline{c}}{c}\right) \right| $$
and
$$|\widehat{u}(\overline{\xi_n})_{k_n}| = |c| \cdot \left|\overline{\sigma_{k_n}(\xi_n)}-\left(\overline{q}-\frac{p \cdot \overline{c}}{c}\right) \right| $$
we must have that $\overline{q}-p \cdot \overline{c}/c \in \mathbb{R}$. Using this observation, we are going to split the proof in some cases.

\begin{itemize} 
	
\item[1)] $\Im(\overline{q}-p) \neq 0$:

In this case we have that $u_1 \in \mathcal{D}'(G) \setminus C^\infty(G)$. In fact, if we suppose that $u_1 \in C^\infty(G)$, then by the observation we made above we must have $\overline{q}-p\cdot 1/1=\overline{q}-p \in \mathbb{R}$, which is a contradiction with our hypothesis. So $u_1 \in \mathcal{D}'(G) \setminus C^\infty(G)$, which implies that $P$ is not (GH).

\item[2)] $\Im(\overline{q}-p) = 0$:

This hypothesis means that $\Im(p)+\Im(q)=0$. This case will be subdivided into other two:

\item[2.1)] $\Im(p) \neq 0$;

For this case we will prove that $u_i \in \mathcal{D}'(G)\setminus C^\infty(G)$. In fact, if $u_i \in C^\infty(G)$, then by the observation above we must have $\overline{q}-p \cdot \overline{i}/i=\overline{q}+p \in \mathbb{R}$, which means that $\Im(q)=\Im(p)$. But since we are assuming that $\Im(p)+\Im(q)=0$, this implies that $\Im(p) =0$, which is a contradiction with our hypothesis. So $u_i \in \mathcal{D}'(G)\setminus C^\infty(G)$ and again $P$ is not (GH).

\item[2.2)] $\Im(p) =0$.

Finally, for this last case observe that since we are assuming $\Im(p)+\Im(q)=0$, the equation $\Im(p)=0$ implies that both $p,q$ are real numbers. We will see that $u_{1+ip} \in \mathcal{D}'(G)\setminus C^\infty(G)$. In fact, if this is not the case, then $\overline{q}-p \cdot (1-ip)/(1+ip)$ should be a real number. But

$$\overline{q}-p \cdot (1-ip)/(1+ip) = q-p \cdot \frac{1-2 p \cdot i -p^2}{1+p^2} =q-\frac{p(1-p^2)}{1+p^2}+\frac{2p^2}{1+p^2} \cdot i,$$
and since $p,q \in \mathbb{R}$, for this number to be real we must have $p=0$, which is a contradiction with our initial hypothesis about $p$. So $u_{1+ip} \in \mathcal{D}'(G)\setminus C^\infty(G)$ and again $P$ is not (GH). This concludes that if the condition (DC) does not hold, then $P$ is not (GH).
\end{itemize}
\end{proof}

\end{document}
